\documentclass[10pt]{article}

\usepackage[utf8]{inputenc}

\usepackage[T1]{fontenc}

\usepackage{amsmath}

\usepackage{amsfonts}

\usepackage{amssymb}

\usepackage[version=4]{mhchem}

\usepackage{stmaryrd}

\usepackage{mathrsfs}

\usepackage{bbold}



\begin{document}

квантованных полей с различными спинами. При этом кванты тех полей, к-рые не зависят от координат $d$-мерного пространства, имеют массу



\[

m \ll m_{p l}=\hbar / l_{p l} c \approx 10^{19} \Gamma \ni \mathrm{B} / c^{2},

\]



а остальные поля - очень тяжелые:



\[

m \geqslant m_{p l}

\]



и поэтому не проявляются в лабораторных экспериментах ( $m_{p l}-$ Планка масса). Наибольший интерес представляет 10 -мерная теория типа К.-К. т. $(d=6)$, к-рая возникает в низкоэнергетич. $E \leqslant m_{p l} c$ пределе более фундаментальной теории нелокальных объектов - суперструн. Нек-рые варианты теории суперструн не содержат ультрафиолетовых расходимостей при специальном выборе группы симметрии $\left(E_{8} \times E_{8}\right.$, к-рая далее нарушается до $\left.E_{6} \times E_{8}\right)$ великого объединения моделей сильного, электромагнитного и слабого взаимодействий. Важно, что эта группа симметрии является удовлетворительной с точки зрения классификации элементарных частиц. А. А. Старобинский.



КАМ-ТЕOPИЯ - раздел теории возмущений и теории малых колебаний, в котором изучаются условно-периодические движения в гамильтоновых системах и родственных им (главным образом обратимых) системах. Название теории образовано по именам ее авторов - А. Н. Колмогорова, В. И. Арнольда и Ю. Мозера (Т. Moser).



Предмет и основные результаты КАМ-т. Центральным результатом КАМ-т. для гамильтоновых уравнений является следующее утверждение.



Т е о р е ма. Типичная гамильтонова система с $n$ степенями свободы обладает для каждого $m, 1 \leqslant m \leqslant n$, семейством $m$-мерных изотропных инвариантных торов, по к-рым происходит условно-периодич. движение с $m$ сильно несоизмеримыми частотами. Объединение всех $m$-мерных инвариантных торов имеет положительную $2 m$-мерную меру.



Слово «типичная» здесь означает, что в функциональном пространстве всех гамильтоновых уравнений, заданных на фиксированном $2 n$-мерном симплектич. многообразии, системы с указанным свойством образуют (в $C^{r}$-топологии для нек-рого достаточно большого $r$ ) непустое открытое множество.



С л е д с т в и е. Гамильтонова система общего вида на гиперповерхностях уровня гамильтониана не эргодична.



Подмногообразие $\mathscr{L}$ симплектич. многообразия $M$ называется и з о т р п п ы м, если ограничение на $\mathscr{L}$ симплектич. структуры равно нулю. Числа $\omega_{1}, \ldots, \omega_{m}$ называются с и л ь но несоизмеримыми типа $(x, \tau)$, если для всех $\boldsymbol{k}=\left(k_{1}, \ldots, k_{m}\right) \in \mathbb{Z}^{m} \backslash\{0\}$



\[

\left|k_{1} \omega_{1}+\ldots+k_{m} \omega_{m}\right| \geqslant x\left(\left|k_{1}\right|+\ldots+\left|k_{m}\right|\right)^{-\tau}

\]



где $x>0$ и $\tau>m-1-$ нек-рые константы, не зависящие от $\boldsymbol{k}$. Эти константы - одни и те же для частот условно-периодич. движения на всех инвариантных торах из данного семейства.



При $m<n$ инвариантные $m$-мерные торы называются не полномер ным и, или маломе рными.



Аналогичные теоремы справедливы для обратимых уравнений, точных симплектич. диффеоморфизмов и обратимых диффеоморфизмов, а также для неавтономных гамильтоновых или обратимых уравнений с периодическими по времени правыми частями. В обратимом случае свойство изотропности торов заменяется на инвариантность относительно обращающей инволюции.



KAM-т. тесно связана с теорией устойчивости движения и теорией адиабатич. инвариантов.



Глобальый и локальный аспекты КАМ-т. Для доказательства существования большого числа инвариантных торов у типичной гамильтоновой или обратимой системы используются следующие два основных подхода.



Гл о бал ь ны Й по д х о д: возмущение интегрируемых систем. Рассматривается «невозмущенная» интегрируемая (полностью или лишь на нек-ром подмногообразии) система, в фазовом пространстве к-рой инвариантные торы дан- ной размерности $m$ составляют гладкую $2 m$-мерную поверхность. Если эта система удовлетворяет нек-рым известным условиям, то при любом достаточно малом и гладком ее возмущении болышинство инвариантных торов не разрушается, а лишь слегка деформируется. Таким образом, окрестность невозмущенной системы в соответствующем функциональном пространстве доставляет искомый пример открытого множества систем, каждая из к-рых имеет много инвариантных торов.



Л о каль н ы й по д х о д: инвариантные торы около положения равновесия, периодич. решения или неподвижной точки. Рассматривается система $\mathfrak{M}$, заданная вблизи критич. элемента $\Gamma$ - положения равновесия (для уравнений), периодич. траектории (для автономных уравнений) или неподвижной точки (для диффеоморфизмов). Если эта система удовлетворяет известным условиям $\mathfrak{A}$, то в каждой окрестности элемента $\Gamma$ она обладает большим числом инвариантных $m$-мерных торов, к-рые при приближении к $\Gamma$ располагаются все гуще. При этом любая система, достаточно близкая к $\mathfrak{R}$, имеет критич. элемент $\Gamma^{\prime}$, близкий к $\Gamma$, и также удовлетворяет условиям $\mathfrak{\cup}$. Таким образом, окрестность системы $\mathfrak{P}$ в соответствуюшем функциональном пространстве доставляет искомый пример открытого множества систем, каждая из к-рых имеет много инвариантных торов.



При локальном подходе роль невозмущенной системы играет нормальная форма около критич. элемента $\Gamma$, а роль величины возмущения - расстояние до Г.



Почти все КАМ-теоремы можно расклассифицировать по следующим признакам:



a) тип симметрии - гамильтоновость или обратимость;



б) тип динамич. системы - автономные дифференциальные уравнения (то есть векторные поля), диффеоморфизмы или неавтономные периодические по времени уравнения;



в) ситуация, в к-рой доказывается сушествование инвариантных торов, - глобальная или локальная;



г) . размерности фазового пространства, многообразия неподвижных точек обращающей инволюции (в обратимом случае) и инвариантных торов;



д) класс гладкости динамич. системы.



Теорема Колмогорова. Значительная часть исследований в КАМ-т. посвящена изучению инвариантных $n$-мерных торов гамильтоновых систем с $n$ степенями свободы, близких к интегрируемым.



Пусть



\[

\dot{I}=-\partial H / \partial \varphi, \dot{\varphi}=\partial H / \partial I

\]



\begin{itemize}

  \item система с гамильтонианом

\end{itemize}



\[

H(I, \varphi, \varepsilon)=H_{0}(I)+\varepsilon H_{1}(I, \varphi, \varepsilon),

\]



где $I=\left(I_{1}, \ldots, I_{n}\right)-$ переменные действия, изменяющиеся в нек-рой (конечной) области $D \subset \mathbb{R}^{n}, \varphi=\left(\varphi_{1}, \ldots, \varphi_{n}\right) \in T^{n}=$ $=(\mathbb{R} / 2 \pi \mathbb{Z})^{n}-$ угловые переменные, а $\varepsilon \in \mathbb{R}-$ малый параметр. Невозмущенная система (с гамильтонианом $H_{0}(I)$ ) интегрируема, и ее фазовое пространство $D \times T^{n}$ расслоено на инвариантные $n$-мерные торы $\{I=$ const $\}$, движение на к-рых условно-периодично с вектором частот $\omega(I)=$ $=\partial H_{0}(I) / \partial I$



Теорема 1 (теорема Колмогорова, или основная теорема КА М-т., см. [1]-[6]). Если гамильтониан $H$ имеет гладкость $C^{r}, r>2 n$, по $I$ и $\varphi$, а невозмущенная система невырождена (det $\|\partial \omega / \partial I\| \neq 0)$ или изоэнергетически невырождена $\operatorname{det}\left\|\begin{array}{cc}\omega / \partial I & \omega \\ \omega & 0\end{array}\right\| \neq 0$, то при достаточно малых $\varepsilon \neq 0$ в фазовом пространстве возмущенной системы содержатся инвариантные $n$-мерные торы (близкие к инвариантным торам невозмущенной системы), по к-рым происходит условно-периодич. движение с сильно несоизмеримыми частотами. Мера дополнения к объединению этих





\end{document}